\chapter{Dinh tuyen dong OSPF voi FRRouting}

\section{Giao thuc OSPF}

OSPF (Open Shortest Path First) la giao thuc dinh tuyen trang thai duong lien ket (Link-State), hoat dong trong mot He tu tri (Autonomous System). Nhung dac diem chinh:

\begin{itemize}
    \item \textbf{Thuat toan Dijkstra (SPF):} Tinh duong di ngan nhat, dam bao khong co vong lap.
    \item \textbf{Hoi tu nhanh:} Khi link down, OSPF phat hien va cap nhat bang dinh tuyen trong vai giay.
    \item \textbf{Ho tro nhieu vung (Areas):} Giam luu luong LSA trong mang lon.
    \item \textbf{Chi phi (Cost):} Metrix dua tren bang thong (100 Mbps / bandwidth), tu dong chon duong toi uu.
\end{itemize}

\section{FRRouting -- Goi phan mem dinh tuyen}

FRRouting (FRR) la giao phan mem dinh tuyen nguon mo chay tren Linux, gom:

\begin{itemize}
    \item \textbf{Zebra daemon:} Quan ly bang dinh tuyen he dieu hanh va dong bo voi kernel.
    \item \textbf{OSPFd:} Daemon chay giao thuc OSPF, gui/nhan LSA va chay thuat toan SPF.
    \item \textbf{Vty shell:} Giao dien CLI tuong tu Cisco IOS de kiem tra trang thai.
\end{itemize}

\section{Cau hinh OSPF cho tung Router}

\subsection{Core Router}

\begin{lstlisting}
! ospfd.conf -- Core Router (router-id 1.1.1.1)
router ospf
  ospf router-id 1.1.1.1
  network 192.168.1.0/24 area 0
  network 203.0.113.0/24 area 0
\end{lstlisting}

\subsection{Distribution Router dist1 va dist2}

\begin{lstlisting}
! dist1: router-id 2.2.2.2
router ospf
  ospf router-id 2.2.2.2
  network 192.168.1.0/29  area 0
  network 172.16.10.0/24  area 1
  redistribute connected

! dist2: router-id 3.3.3.3
router ospf
  ospf router-id 3.3.3.3
  network 192.168.1.4/30  area 0
  network 172.16.20.0/24  area 1
  redistribute connected
\end{lstlisting}

\subsection{DMZ Router}

\begin{lstlisting}
! dmz_r: router-id 4.4.4.4
router ospf
  ospf router-id 4.4.4.4
  network 192.168.1.8/30  area 0
  network 10.10.10.0/24   area 2
  redistribute connected
\end{lstlisting}

\section{Thiet ke OSPF Areas}

\begin{table}[H]
\centering
\caption{Phan chia OSPF Areas trong he thong Campus Network}
\vspace{0.3cm}
\begin{tabularx}{\textwidth}{|c|c|>{\raggedright\arraybackslash}X|>{\raggedright\arraybackslash}X|}
    \hline
    \textbf{Area} & \textbf{Loai} & \textbf{Subnet} & \textbf{Router} \\ \hline
    Area 0 (Backbone) & Normal & 192.168.1.0/24, 203.0.113.0/24 & Core, dist1, dist2, dmz\_r \\ \hline
    Area 1 & Normal & 172.16.10.0/24, 172.16.20.0/24 & dist1, dist2 \\ \hline
    Area 2 & Stub & 10.10.10.0/24 & dmz\_r \\ \hline
\end{tabularx}
\label{tab:ospf_areas}
\end{table}

\section{Kiem tra OSPF sau khi trien khai}

Sau khi khoi dong FRR tren tat ca cac Router, kiem tra trang thai OSPF bang vtysh:

\begin{lstlisting}
# Kiem tra Neighbor
show ip ospf neighbor
# -> dist1, dist2, dmz_r phai o trang thai Full/DR hoac Full/BDR

# Kiem tra Bang dinh tuyen OSPF
show ip route ospf
# O  172.16.10.0/24 [110/20] via 192.168.1.2
# O  172.16.20.0/24 [110/20] via 192.168.1.6
# O  10.10.10.0/24  [110/20] via 192.168.1.10
\end{lstlisting}

\textbf{Ket qua hoi tu khi su co:} Khi tat ket noi Core -- dist1, OSPF tu dong hoi tu qua dist2 trong khoang 12~--~15 giay, khong can chinh sua thu cong.
